\documentclass[11pt,a4paper,oneside,openright]{book}
\usepackage[spanish,activeacute,es-noshorthands]{babel}
\usepackage[utf8]{inputenc}
\usepackage[T1]{fontenc}
\usepackage{lmodern}
\usepackage{amsmath, amssymb, amsthm, mathtools}
\usepackage{graphicx}
\usepackage[dvipsnames]{xcolor}
\usepackage{tikz}
\usetikzlibrary{circuits.ee.IEC, calc, arrows.meta, decorations.pathreplacing}
\usepackage{listings}
\usepackage{geometry}
\geometry{margin=2.5cm, headheight=14pt}
\usepackage{microtype}
\usepackage{booktabs}
\usepackage{enumitem}
\usepackage{fancyhdr}
\usepackage{tcolorbox}
\tcbuselibrary{skins, breakable}
\usepackage{caption}
\captionsetup{font=small, labelfont=bf}
\usepackage{hyperref}
\hypersetup{
  colorlinks=true,
  linkcolor=NavyBlue,
  citecolor=ForestGreen,
  urlcolor=BrickRed,
  pdftitle={Manual de electronica practica para CubeSat-EdgeAI},
  pdfauthor={INTISAT}
}

\newtheorem{definicion}{Definicion}[chapter]
\newtheorem{teorema}[definicion]{Teorema}
\newtheorem{proposicion}[definicion]{Proposicion}
\newtheorem*{nota}{Nota}

\definecolor{codebg}{RGB}{248, 249, 250}
\definecolor{codecomment}{RGB}{106, 153, 85}
\definecolor{codekey}{RGB}{86, 156, 214}
\definecolor{codestr}{RGB}{206, 145, 120}

\lstdefinestyle{python}{
  language=Python,
  basicstyle=\ttfamily\footnotesize,
  keywordstyle=\color{codekey}\bfseries,
  commentstyle=\color{codecomment}\itshape,
  stringstyle=\color{codestr},
  numberstyle=\tiny\color{gray},
  numbers=left, numbersep=8pt,
  frame=single, rulecolor=\color{gray!30},
  framesep=4pt,
  backgroundcolor=\color{codebg},
  showstringspaces=false, breaklines=true, breakatwhitespace=true,
  tabsize=2, captionpos=b,
  literate={a}{{a}}1 {e}{{e}}1 {i}{{i}}1 {o}{{o}}1 {u}{{u}}1 {n}{{n}}1
}
\lstset{style=python}

\newtcolorbox{intuibox}[1][]{
  colback=green!5, colframe=green!40!black,
  fonttitle=\bfseries, title=Intuicion,
  breakable, sharp corners=south, #1
}
\newtcolorbox{calculo}[1][]{
  colback=gray!4, colframe=gray!55,
  fonttitle=\bfseries, title=Calculo,
  breakable, sharp corners=south, #1
}
\newtcolorbox{payload}[1][]{
  colback=blue!4, colframe=blue!50!black,
  fonttitle=\bfseries, title=Conexion con el payload,
  breakable, sharp corners=south, #1
}
\newtcolorbox{cuidado}[1][]{
  colback=red!5, colframe=red!50!black,
  fonttitle=\bfseries, title=Cuidado,
  breakable, sharp corners=south, #1
}
\newtcolorbox{ejercicio}[1][]{
  colback=orange!5, colframe=orange!60!black,
  fonttitle=\bfseries, title=Ejercicios,
  breakable, sharp corners=south, #1
}

\pagestyle{fancy}
\fancyhf{}
\fancyhead[LE,RO]{\thepage}
\fancyhead[RE]{\nouppercase{\leftmark}}
\fancyhead[LO]{\nouppercase{\rightmark}}

\title{\Huge\bfseries Manual de electronica practica\\[0.4em]
\Large De los electrones al CubeSat\\[0.6em]
\large Atado al proyecto CubeSat-EdgeAI-Payload}
\author{Material acompañante de \texttt{INTISAT}}
\date{\today}

\begin{document}

\frontmatter
\maketitle

\chapter*{Prefacio}
\addcontentsline{toc}{chapter}{Prefacio}

Este manual es para alguien que se siente \emph{perdido} cuando aparecen
terminos como I2C, SPI, INA219, MOSFET, regulador LDO, ground loop, etc.
La idea es construir el lenguaje desde lo mas basico (que es un electron,
que es un voltaje) hasta poder mantener una conversacion tecnica sobre
diseño de carrier board para un satelite.

\section*{Como leer este libro}

El libro tiene 7 partes que avanzan en complejidad. Las primeras 3 son
\emph{base obligatoria}; las siguientes son \emph{aplicaciones} a tu
proyecto.

\begin{description}[leftmargin=2em, style=nextline]
  \item[Parte I --- Fundamentos electricos.] Que es voltaje, corriente,
        resistencia. Ley de Ohm. Componentes pasivos.
  \item[Parte II --- Componentes activos.] Diodos, transistores, op-amps.
        Reguladores. La diferencia entre pasivo y activo.
  \item[Parte III --- Electronica digital.] 0 y 1, compuertas, microcontroladores
        y SoCs.
  \item[Parte IV --- Comunicacion entre chips.] I2C, SPI, UART. Niveles
        logicos. Level shifters.
  \item[Parte V --- Sensores y actuadores.] CMOS, INA219, OLED, camera CSI.
  \item[Parte VI --- Sistemas de potencia.] Bateries, reguladores, presupuesto
        energetico, EPS satelital.
  \item[Parte VII --- Aplicacion al CubeSat.] El stack del payload INTISAT,
        el carrier board, los tests pre-vuelo.
\end{description}

Cada capitulo tiene:
\begin{enumerate}[label=\arabic*., leftmargin=2em]
  \item \textbf{Por que importa} --- el problema concreto que resuelve.
  \item \textbf{Intuicion} --- la analogia antes de la formula (caja verde).
  \item \textbf{Matematica} --- definiciones formales y ecuaciones (caja gris).
  \item \textbf{Conexion con el payload} --- donde aparece en el INTISAT
        (caja azul).
  \item \textbf{Cuidado} --- errores comunes (caja roja).
  \item \textbf{Ejercicios} --- experimentos sobre tu hardware (caja naranja).
\end{enumerate}

\section*{Prerequisitos}

Solo se asume:
\begin{itemize}[leftmargin=2em]
  \item Algebra basica (resolver $V = I R$).
  \item Concepto de funciones (sin necesidad de derivar/integrar fuera de
        algunos casos puntuales).
  \item Computadora con Python para simular y verificar.
\end{itemize}

\textbf{No se asume} conocimiento previo de electronica. Empezamos de cero.

\tableofcontents

\mainmatter

% =====================================================================
\part{Fundamentos electricos}

\chapter{Magnitudes y unidades}

\section{Por que importa}

Antes de hablar de cualquier circuito, hay que hablar el mismo idioma. Cada
medida que vas a leer en un datasheet (corriente de la Pi, voltaje de la
bateria, capacidad de un capacitor) usa una unidad estandar del SI. Te paso
las 7 que vas a usar el 95\% del tiempo.

\section{Las 7 magnitudes esenciales}

\begin{center}
\begin{tabular}{lllll}
\toprule
\textbf{Magnitud} & \textbf{Simbolo} & \textbf{Unidad SI} & \textbf{Sigla} & \textbf{Lo que mide} \\
\midrule
Voltaje (tension)   & $V, U$ & volt     & V    & "presion electrica" \\
Corriente            & $I$    & ampere   & A    & cantidad de carga por segundo \\
Resistencia          & $R$    & ohm      & $\Omega$ & oposicion al paso de corriente \\
Potencia             & $P$    & watt     & W    & energia por unidad de tiempo \\
Energia              & $E$    & joule    & J    & "trabajo total" \\
Capacitancia         & $C$    & farad    & F    & cuanta carga acumula al voltaje \\
Inductancia          & $L$    & henry    & H    & cuanta energia magnetica almacena \\
\bottomrule
\end{tabular}
\end{center}

\section{Voltaje --- la "presion"}

\begin{intuibox}
Pensa en un sistema de cañerias con agua. El \textbf{voltaje} es como la
\emph{presion del agua}: mide cuanta fuerza tiene el sistema para empujar
el flujo. Un punto a 5 V respecto a tierra (GND) tiene 5 V de "presion
electrica" mas que el GND.
\end{intuibox}

\begin{definicion}[Voltaje]
El voltaje entre dos puntos $A$ y $B$ es la diferencia de potencial
electrico entre ellos:
\begin{equation}
V_{AB} = V_A - V_B
\end{equation}
Un voltaje absoluto no existe: siempre se mide \emph{respecto a otro punto}
(generalmente GND, que arbitrariamente se llama 0 V).
\end{definicion}

\textbf{Voltajes que vas a ver}:

\begin{center}
\begin{tabular}{lc}
\toprule
\textbf{Sistema} & \textbf{Voltaje tipico} \\
\midrule
Bateria CR2032 (reloj de pulsera)   & 3.0 V \\
USB                                  & 5.0 V \\
Pi 5 (input)                         & 5.0 V \\
GPIO de la Pi                        & 3.3 V \\
Bateria de auto                      & 12 V \\
Red electrica (Argentina)            & 220 V (AC) \\
Linea de alta tension                & 13 200 V \\
\bottomrule
\end{tabular}
\end{center}

\section{Corriente --- el flujo}

\begin{intuibox}
Si el voltaje es la \emph{presion}, la corriente es el \emph{caudal de
agua} que efectivamente pasa. Mide cuantos electrones por segundo cruzan
una seccion del cable.
\end{intuibox}

\begin{definicion}[Corriente]
Corriente electrica es la cantidad de carga $Q$ que pasa por un punto en
un tiempo $t$:
\begin{equation}
I = \frac{Q}{t}
\end{equation}
Unidad: ampere (A) $=$ coulomb por segundo (C/s). Un coulomb son
$6.24 \times 10^{18}$ electrones.
\end{definicion}

\textbf{Corrientes que vas a ver}:

\begin{center}
\begin{tabular}{lc}
\toprule
\textbf{Caso} & \textbf{Corriente tipica} \\
\midrule
LED encendido                & 10--20 mA \\
Sensor OV5647 activo          & 50--200 mA \\
Pi 5 idle                     & 600 mA \\
Pi 5 carga 100\%              & 1.4--1.8 A \\
Boot transient Pi 5           & hasta 2.0 A \\
Hervidor electrico            & 8--10 A \\
Linea principal de la casa    & 30--60 A \\
\bottomrule
\end{tabular}
\end{center}

\subsection{Convencion de signo}

\textbf{Corriente convencional} fluye de + a $-$ (del punto de mayor a
menor voltaje). Los electrones fisicamente se mueven al reves (de $-$ a
+) pero se acordo en el siglo XIX usar la convencion contraria. Hoy todo
el mundo trabaja con la convencional.

\section{Resistencia --- la oposicion}

\begin{intuibox}
Volviendo a la cañeria: una \textbf{resistencia} es como una manguera
mas angosta que dificulta el paso del agua. A mas resistencia, menos caudal
para la misma presion.
\end{intuibox}

\begin{definicion}[Resistencia]
Es la oposicion al paso de corriente. Su unidad es el ohm ($\Omega$):
1 $\Omega$ es la resistencia que deja pasar 1 A cuando se aplica 1 V.
\end{definicion}

\section{La Ley de Ohm}

\begin{teorema}[Ley de Ohm]
Para una resistencia ideal, la corriente es proporcional al voltaje:
\begin{equation}
\boxed{V = I \cdot R}
\label{eq:ohm}
\end{equation}
\end{teorema}

\textbf{Como leerla}:
\begin{itemize}[leftmargin=2em]
  \item Si fijas $R$, doblar el voltaje dobla la corriente.
  \item Si fijas $V$, doblar la resistencia parte la corriente al medio.
\end{itemize}

\begin{calculo}
\textbf{Ejemplo 1}: una LED roja necesita $20$ mA y cae $2$ V. La conectas
a $5$ V. ¿Que resistencia poner en serie?

La resistencia tiene que caer $5 - 2 = 3$ V con $0.020$ A pasando:
\[
R = \frac{V}{I} = \frac{3}{0.020} = 150\;\Omega
\]
\end{calculo}

\begin{calculo}
\textbf{Ejemplo 2}: tu Pi 5 idle consume $600$ mA a $5$ V. ¿Cual es su
resistencia equivalente?
\[
R = \frac{V}{I} = \frac{5}{0.6} = 8.3\;\Omega
\]
La Pi se "ve" como una resistencia de $8.3\;\Omega$ desde la EPS.
\end{calculo}

\section{Potencia --- la energia por segundo}

\begin{definicion}[Potencia]
\begin{equation}
\boxed{P = V \cdot I}
\label{eq:power}
\end{equation}
Unidad: watt (W) $=$ volt $\times$ ampere.
\end{definicion}

Combinado con Ohm:
\begin{align}
P &= V \cdot I = I^2 \cdot R = \frac{V^2}{R}
\end{align}

\textbf{Ejemplos de potencia}:
\begin{center}
\begin{tabular}{lc}
\toprule
LED indicador & 0.04 W \\
Pi 5 idle & 3 W \\
Pi 5 a tope & 8 W \\
Lampara LED de habitacion & 9 W \\
Notebook & 50 W \\
Motor de auto electrico & 100 000 W (100 kW) \\
\bottomrule
\end{tabular}
\end{center}

\section{Energia}

La energia es potencia $\times$ tiempo:
\begin{equation}
E = P \cdot t
\end{equation}
Unidad: joule (J) $=$ watt $\times$ segundo. Tambien se mide en
\textbf{watt-hora} (Wh): $1$ Wh $= 3600$ J.

\textbf{Ejemplo}: Pi 5 corriendo a $7.5$ W durante $30$ s consume:
\[
E = 7.5 \times 30 = 225\;\text{J} = \frac{225}{3600} \approx 0.063\;\text{Wh}
\]

\begin{payload}
Tu cuadro 4 del reporte energetico decia: \emph{286 J en modo FPM, 23 J en
modo single}. Esos numeros vienen de aplicar $E = P \cdot t$ con la
estimacion de $6.5$ W de potencia tipica y los tiempos medidos. Como NO
mediste $P$ con INA219, asumiste $6.5$ W de datasheet. Esa es la limitacion.
\end{payload}

\begin{ejercicio}
\begin{enumerate}[leftmargin=2em]
  \item ¿Cuanta corriente consume un LED a $20$ mA si lo conectas a una
        bateria de $9$ V con una resistencia de $470\;\Omega$? (sugerencia:
        primero calcula el voltaje sobre el LED).
  \item Tu Pi 5 corre $10$ scans en modo single ($23$ J cada uno). Una
        bateria de $5000$ mAh a $5$ V tiene $\sim 90\,000$ J. ¿Cuantos scans
        soporta?
  \item Si la EPS del INTISAT da $5$ V $\times$ $6$ A, ¿cual es la potencia
        maxima del rail \texttt{5V\_PAYLOAD}? ¿Cuanto margen tenes vs el
        pico de Pi 5 a $8$ W?
\end{enumerate}
\end{ejercicio}

\chapter{Componentes pasivos}

\section{Resistor}

El resistor es el componente mas simple: un trozo de material resistivo
encapsulado. Se identifica por:

\begin{itemize}[leftmargin=2em]
  \item \textbf{Valor}: en ohms, marcado con codigo de colores o numero.
  \item \textbf{Tolerancia}: $\pm 1\%$, $\pm 5\%$, $\pm 10\%$.
  \item \textbf{Potencia maxima}: 1/8 W, 1/4 W, 1/2 W, 1 W, etc.
\end{itemize}

\begin{cuidado}
Si un resistor de $1/4$ W disipa $1$ W, se quema literalmente. Verificar
siempre $P_{\text{esperado}} = I^2 R < P_{\text{nominal}}$.
\end{cuidado}

\subsection{En serie y en paralelo}

\textbf{Serie} (uno detras del otro): las resistencias se suman.
\begin{equation}
R_{\text{eq}} = R_1 + R_2 + \dots + R_n
\end{equation}

\textbf{Paralelo} (en lineas separadas): los reciprocos se suman.
\begin{equation}
\frac{1}{R_{\text{eq}}} = \frac{1}{R_1} + \frac{1}{R_2} + \dots + \frac{1}{R_n}
\end{equation}

Para 2 en paralelo: $R_{\text{eq}} = \frac{R_1 R_2}{R_1 + R_2}$.

\begin{intuibox}
\textbf{Serie} es como poner mas filtros en linea: cada uno suma su
oposicion. \textbf{Paralelo} es abrir mas caminos: aunque cada uno tenga
oposicion, hay mas vias y la resistencia total baja.
\end{intuibox}

\section{Capacitor}

\begin{definicion}[Capacitor]
Es un componente que almacena energia en forma de campo electrico. Su
ecuacion es:
\begin{equation}
Q = C \cdot V
\end{equation}
donde $Q$ es la carga acumulada y $C$ la capacitancia (F).
\end{definicion}

\textbf{Para entenderlo}: cuando aplicas un voltaje, el capacitor "se llena"
de carga; cuando lo desconectas, mantiene esa carga (y por ende voltaje)
hasta que algo lo descargue.

\subsection{Tipos comunes}

\begin{center}
\begin{tabular}{llll}
\toprule
\textbf{Tipo} & \textbf{Capacitancia} & \textbf{Voltaje max} & \textbf{Uso tipico} \\
\midrule
Ceramico               & pF a $\mu$F   & 50--100 V & filtros, decoupling \\
Electrolitico aluminio & $\mu$F a mF   & 6.3--400 V & filtros de fuente \\
Tantalio               & $\mu$F a 100$\mu$F & 6.3--50 V & decoupling preciso \\
Cermet                 & pF a nF       & alta       & RF \\
\bottomrule
\end{tabular}
\end{center}

\subsection{Comportamiento dinamico}

A diferencia del resistor (que cumple Ohm de forma instantanea), el
capacitor reacciona al cambio:
\begin{equation}
I = C \frac{dV}{dt}
\end{equation}

\begin{intuibox}
Un capacitor se opone a los \emph{cambios bruscos} de voltaje. Por eso se
usa para filtrar ruido: el ruido es un cambio rapido y el capacitor lo
absorbe.
\end{intuibox}

\subsection{En serie y paralelo}

\textbf{Serie}: capacitancias en reciproco.
\textbf{Paralelo}: se suman directamente (lo opuesto a las resistencias).

\section{Inductor}

\begin{definicion}[Inductor]
Almacena energia en campo magnetico. Su ecuacion:
\begin{equation}
V = L \frac{dI}{dt}
\end{equation}
Se opone a los cambios de \emph{corriente}.
\end{definicion}

Aparece tipicamente en:
\begin{itemize}[leftmargin=2em]
  \item Reguladores switching (buck, boost, SEPIC).
  \item Filtros de RF.
  \item Antenas.
\end{itemize}

\begin{payload}
En tu carrier board para CM5 vas a poner capacitores de \textbf{decoupling}
de 100 nF cerca de cada chip. Filtran ruido del rail 3.3 V. Sin ellos,
los chips fallan intermitentemente. Es la regla \#1 de PCB design.
\end{payload}

\chapter{Leyes de Kirchhoff}

\section{Por que importa}

Ohm sirve para una sola resistencia. Para circuitos con varios componentes
necesitas las dos leyes de Kirchhoff. Son la base de todo analisis de
circuitos.

\section{KCL --- Ley de corrientes}

\begin{teorema}[Kirchhoff Current Law]
La suma de corrientes que entran a un nodo es igual a la suma de las que
salen:
\begin{equation}
\sum I_{\text{entrada}} = \sum I_{\text{salida}}
\end{equation}
\end{teorema}

\begin{intuibox}
Es la conservacion de la carga: lo que entra a un punto tiene que salir
(no se acumula). Como en una cañeria con un nodo, el agua que llega es la
que sale por las ramas.
\end{intuibox}

\section{KVL --- Ley de voltajes}

\begin{teorema}[Kirchhoff Voltage Law]
La suma de voltajes en un lazo cerrado es cero:
\begin{equation}
\sum V_{\text{lazo}} = 0
\end{equation}
\end{teorema}

\begin{intuibox}
Si recorres un circuito cerrado y vuelves al mismo punto, la suma de
"subidas" y "bajadas" de voltaje tiene que dar cero.
\end{intuibox}

\section{Aplicacion: divisor de voltaje}

Dado dos resistencias en serie $R_1$, $R_2$ con voltaje total $V_{\text{in}}$:
\begin{equation}
V_{\text{out}} = V_{\text{in}} \cdot \frac{R_2}{R_1 + R_2}
\end{equation}

Es la formula mas usada en electronica analogica.

\begin{ejercicio}
\begin{enumerate}[leftmargin=2em]
  \item Un divisor con $R_1 = 10$ k$\Omega$ y $R_2 = 5$ k$\Omega$ alimentado
        con 5 V. ¿Cuanto sale?
  \item ¿Que pasa si conectas una carga de $1$ k$\Omega$ a la salida?
        (sugerencia: la carga queda en paralelo con $R_2$).
\end{enumerate}
\end{ejercicio}

% =====================================================================
\part{Componentes activos}

\chapter{Diodos}

\section{Que es}

Un diodo es un componente de \textbf{2 terminales} que solo deja pasar
corriente en una direccion. Es como una valvula check de plomeria: si la
corriente trata de ir al reves, la bloquea.

\subsection{Curva caracteristica}

Se polariza directo: $V_{\text{anodo}} > V_{\text{catodo}} + V_F$
($V_F \approx 0.7$ V para Si, $\approx 0.3$ V para Schottky).

\subsection{Aplicaciones}

\begin{itemize}[leftmargin=2em]
  \item \textbf{Rectificacion}: convertir AC a DC.
  \item \textbf{Proteccion de polaridad inversa}: si conectas la fuente al
        reves, el diodo bloquea.
  \item \textbf{Diodo zener}: estabiliza voltaje de referencia.
  \item \textbf{LED}: emite luz cuando hay corriente directa.
\end{itemize}

\begin{payload}
En tu carrier board CM5 conviene poner un diodo Schottky en serie con el
rail 5V de la EPS, asi si alguien lo conecta al reves, no quemas la Pi.
Costo: $0.10$ USD. Vale ese seguro.
\end{payload}

\chapter{Transistores BJT y MOSFET}

\section{Concepto general}

Un transistor es un \textbf{interruptor electronico controlable}: una pequeña
señal en un terminal controla un flujo grande entre otros dos. Es la base
de la electronica moderna.

\section{BJT (Bipolar Junction Transistor)}

3 terminales: base (B), colector (C), emisor (E). Una pequeña corriente
en la base permite una corriente mucho mayor entre C y E:
\begin{equation}
I_C = \beta \cdot I_B \quad (\beta \approx 100)
\end{equation}

Tipos: NPN (interruptor cierra con base alta) y PNP (cierra con base baja).

\section{MOSFET}

3 terminales: gate (G), drain (D), source (S). Un voltaje en G controla la
corriente D-S, sin consumir corriente en G.

\textbf{Por que importa}: la mayoria de la electronica digital moderna usa
MOSFETs. Un microcontrolador tiene millones de ellos en el chip.

\begin{payload}
Si tu carrier board necesita encender/apagar el rail OLED bajo control de
software, lo haces con un MOSFET P-channel: GPIO de la Pi controla el gate,
el MOSFET deja pasar 5V al OLED. La Pi puede asi cortar el OLED para
ahorrar energia en SAFE\_MODE.
\end{payload}

\chapter{Reguladores de tension}

\section{Que problema resuelve}

Tu fuente da $5$ V pero la Pi internamente necesita $0.85$ V para el CPU,
$1.1$ V para la RAM, $1.8$ V para los buses, $3.3$ V para GPIO. Todos esos
voltajes se generan a partir de los $5$ V externos con \textbf{reguladores}.

\section{LDO (Low Dropout)}

Es un regulador lineal: convierte un voltaje mayor en uno menor disipando
la diferencia como \textbf{calor}.

\textbf{Eficiencia}:
\begin{equation}
\eta = \frac{V_{\text{out}}}{V_{\text{in}}}
\end{equation}

Si convertis $5$ V a $3.3$ V con $1$ A:
\begin{itemize}[leftmargin=2em]
  \item Disipa $(5 - 3.3) \times 1 = 1.7$ W como calor.
  \item Eficiencia: $66\%$.
\end{itemize}

\textbf{Pros}: simple, barato, bajo ruido.
\textbf{Contras}: ineficiente, calienta.

\section{Switching (buck, boost)}

Usa inductor + capacitor para convertir voltaje con $> 90\%$ de eficiencia.
Mas complejo, mas ruidoso, pero mucho mas eficiente.

\textbf{Tipos}:
\begin{itemize}[leftmargin=2em]
  \item \textbf{Buck}: baja voltaje (5V $\to$ 3.3V).
  \item \textbf{Boost}: sube voltaje (3.7V $\to$ 12V).
  \item \textbf{Buck-boost}: ambos.
  \item \textbf{SEPIC}: aislado.
\end{itemize}

\begin{payload}
La PMIC interna de la Pi 5 (chip Renesas) hace todo esto en un solo IC:
recibe 5V externos y genera 0.85V para CPU, 1.1V RAM, 1.8V IO, 3.3V GPIO.
Si lees el datasheet del Pi 5 vas a ver el bloque de la PMIC. Es transparente
para vos como usuario.
\end{payload}

% =====================================================================
\part{Electronica digital}

\chapter{Sistemas de numeracion}

\section{Binario}

Solo dos digitos: 0 y 1. Cada digito representa una potencia de 2.
\[
1011_2 = 1\cdot 8 + 0\cdot 4 + 1\cdot 2 + 1\cdot 1 = 11_{10}
\]

\section{Hexadecimal}

16 digitos: 0--9, A--F. Cada digito hex representa 4 bits.
\[
0xA5 = A\cdot 16 + 5 = 10\cdot 16 + 5 = 165_{10} = 10100101_2
\]

\subsection{Por que importa}

Las direcciones I2C, los registros, los comandos de tu daemon estan en hex.
Tu I2C slave responde en \texttt{0x42}, lo que es:
\[
0x42 = 4 \cdot 16 + 2 = 66_{10} = 0100\,0010_2
\]

\section{Bytes y multiplos}

\begin{center}
\begin{tabular}{lll}
\toprule
\textbf{Unidad} & \textbf{Bytes} & \textbf{Ejemplo} \\
\midrule
1 byte (B)         & 8 bits       & un caracter ASCII \\
1 kB ($2^{10}$ B)  & 1024 B        & un mensaje de texto largo \\
1 MB ($2^{20}$ B)  & 1.05 M        & una imagen JPEG mediana \\
1 GB ($2^{30}$ B)  & 1.07 G        & una pelicula comprimida \\
\bottomrule
\end{tabular}
\end{center}

\chapter{Algebra de Boole y compuertas logicas}

\section{Compuertas basicas}

\begin{center}
\begin{tabular}{lll}
\toprule
\textbf{Compuerta} & \textbf{Operacion} & \textbf{Tabla de verdad} \\
\midrule
AND   & $A \wedge B$           & 1 solo si ambas son 1 \\
OR    & $A \vee B$             & 1 si al menos una es 1 \\
NOT   & $\bar{A}$              & invierte \\
XOR   & $A \oplus B$           & 1 si son diferentes \\
NAND  & $\overline{A \wedge B}$ & complemento de AND \\
NOR   & $\overline{A \vee B}$  & complemento de OR \\
\bottomrule
\end{tabular}
\end{center}

Toda la electronica digital se construye combinando estas. Un microprocesador
con miles de millones de transistores no es mas que estas compuertas
ordenadas.

\chapter{Microcontroladores y SoCs}

\section{Microcontrolador}

Un \textbf{microcontrolador} es un chip que integra: CPU + RAM + flash +
perifericos (GPIO, I2C, SPI, UART, ADC, timers). Ejemplos: Arduino UNO
(ATmega328), STM32, ESP32.

\section{SoC (System on Chip)}

Un \textbf{SoC} es un chip mas potente: CPU multi-core + GPU + ISP + DMA +
controladores de memoria + buses. Ejemplos: el BCM2712 del Pi 5, los
Apple M1/M2, los Qualcomm Snapdragon.

\textbf{Diferencia clave}: el microcontrolador corre firmware; el SoC corre
sistema operativo (Linux).

\begin{payload}
Tu Pi 5 tiene el SoC \textbf{BCM2712}: 4 nucleos ARM Cortex-A76 a 2.4 GHz,
GPU VideoCore VII, RAM LPDDR4X de 4 GB. Lo controla un Linux completo.
Si quisieras hacer un controlador de motor o un sensor simple, usarias un
microcontrolador (ej. ATmega) que es 100x mas barato pero solo soporta
firmware bare-metal.
\end{payload}

% =====================================================================
\part{Comunicacion entre chips}

\chapter{I2C}

\section{Concepto}

\textbf{I2C} = Inter-Integrated Circuit. Bus de 2 cables (mas GND) para
conectar varios chips en serie. Inventado por Philips en 1982.

\section{Cables}

\begin{itemize}[leftmargin=2em]
  \item \textbf{SDA} (Serial Data): los datos.
  \item \textbf{SCL} (Serial Clock): el reloj.
  \item GND: tierra comun.
\end{itemize}

Ambas lineas son \textbf{open-drain}: el chip solo puede tirar la linea a
0 V; nadie puede ponerla a $V_{cc}$. Por eso se necesitan \textbf{pull-up
resistors} (4.7 k$\Omega$ tipico) que mantienen la linea en alto cuando
nadie la baja.

\section{Direcciones}

Cada slave tiene una direccion de 7 bits ($0$--$127$). El master inicia
la comunicacion con: \texttt{[start] [address+R/W] [datos] [stop]}.

\begin{payload}
Tu I2C slave (la Pi) responde en \textbf{0x42}. El INA219 responde en
\textbf{0x40}. El OLED en \textbf{0x3C}. Pueden compartir el mismo bus
sin conflicto. \texttt{i2cdetect -y 1} te muestra todos los slaves
presentes.
\end{payload}

\section{Velocidades estandar}

\begin{center}
\begin{tabular}{lll}
\toprule
\textbf{Modo} & \textbf{Velocidad} & \textbf{Cuando} \\
\midrule
Standard       & 100 kHz          & default, casi todos los chips \\
Fast           & 400 kHz          & sensores rapidos \\
Fast plus      & 1 MHz            & raro \\
High speed     & 3.4 MHz          & muy raro \\
Ultra          & 5 MHz            & muy moderno \\
\bottomrule
\end{tabular}
\end{center}

\chapter{SPI y UART}

\section{SPI (Serial Peripheral Interface)}

4 cables: SCLK, MOSI, MISO, CS. Mas rapido que I2C (hasta $50$ MHz),
full-duplex (lee y escribe a la vez).

\section{UART (Universal Asynchronous)}

2 cables: TX, RX. Sin reloj compartido, depende de baud rate sincronizado
en ambos extremos. Tipicamente 9600, 115200 bps.

\chapter{Niveles logicos y level shifters}

\section{El problema}

Distintos chips usan distintos voltajes para "alto" (1) y "bajo" (0):
\begin{itemize}[leftmargin=2em]
  \item Pi 5 GPIO: $3.3$ V
  \item Arduino UNO: $5$ V
  \item ESP32 logico: $3.3$ V
  \item Logica industrial CMOS: $5$ V o $24$ V
\end{itemize}

Conectar directo dañas chips.

\section{Solucion: level shifter}

Un \textbf{level shifter} es un IC que traduce voltajes en ambos sentidos
sin conectarlos electricamente. Ejemplos: TXS0108E (8 canales),
74LVC245 (octal), BSS138 (MOSFET simple).

\begin{payload}
Tu carrier board para CM5 va a tener un \textbf{TXS0108E} o similar entre
el I2C de la Pi (3.3 V) y el bus PC-104 del INTISAT (que podria estar a
3.3 V o 5 V segun el OBC). Sin el level shifter, riesgo de daño cuando
el OBC tira la linea a $V_{cc}$ que no aplica.
\end{payload}

% =====================================================================
\part{Sensores y actuadores}

\chapter{Sensor de imagen CMOS (OV5647)}

\section{Como funciona}

Un sensor CMOS tiene una grilla de \textbf{fotodiodos} (uno por pixel).
Cuando un foton llega, libera un electron. La cantidad de electrones
acumulados en un tiempo de exposicion es proporcional a la intensidad
luminosa de ese pixel. Despues, un convertidor ADC mide cada pixel y
saca el valor.

\section{Filtro de Bayer}

El sensor es monocromatico por defecto. Para color, se le pone una grilla
de filtros R, G, B alternados (patron Bayer). Cada pixel ve solo un color;
el ISP interpola los otros dos.

\begin{payload}
Tu OV5647 es de $5$ MP ($2592 \times 1944$) con pixeles de $1.4\;\mu$m. Sin
lente, captura proyecciones directas de la luz. El sensor por si solo da
RAW10 (10 bits crudos por pixel). El ISP del Pi 5 lo convierte a RGB888.
\end{payload}

\chapter{INA219 --- sensor de corriente}

\section{Como mide}

Como vimos: ley de Ohm $I = V_{\text{shunt}}/R_{\text{shunt}}$. El INA219
mide el voltaje sobre un shunt de $0.1\;\Omega$ y deduce la corriente.

\begin{payload}
Tu plan es usarlo para medir el consumo real de la Pi durante cada fase
del daemon. Conectalo en serie con el rail $5$ V de la EPS, lee el bus
I2C secundario, y obtenes V, I, P en tiempo real.
\end{payload}

\chapter{OLED SSD1351 / SSD1306}

\section{Como funciona un OLED}

Un OLED (Organic LED) tiene una matriz de pixeles donde cada uno es un
LED organico. Cuando aplicas un voltaje, el pixel emite luz directa
(no necesita backlight como un LCD).

Cada pixel se controla por una tension entre $0$ V (apagado) y $3.3$ V
(brillo maximo). Un controlador (SSD1351 / SSD1306) maneja la matriz por
SPI o I2C.

\begin{payload}
Tu OLED es la fuente de iluminacion del payload. Lo manejas por SPI desde
la Pi a $1$ MHz. En cada captura encendes un patron especifico (5x5
para FPM, full white para single, etc.). Es un actuador opto-electronico
programable.
\end{payload}

% =====================================================================
\part{Sistemas de potencia}

\chapter{Bateria de Litio}

\section{Quimica}

Una bateria Li-ion almacena energia en forma de iones de litio que
migran entre el catodo (LiCoO$_2$) y el anodo (grafito) durante la carga
y descarga.

\section{Voltajes}

\begin{itemize}[leftmargin=2em]
  \item Cargada: $4.2$ V
  \item Nominal: $3.7$ V
  \item Descargada: $3.0$ V (no bajar de aqui o se daña)
\end{itemize}

\section{Capacidad}

Se mide en \textbf{mAh} (mili-amperes-hora) o Wh (watt-hora):
\[
\text{Wh} = \text{V} \times \text{Ah}
\]

Una bateria de $5000$ mAh @ $3.7$ V tiene $5\,\text{Ah} \times 3.7\,\text{V} = 18.5$ Wh.

\chapter{Reguladores switching y LDO}

(Capitulo 8 ya cubrio la teoria.)

\chapter{Presupuesto energetico}

\section{Como se calcula}

Se enumeran todos los modos de operacion del sistema, sus duraciones
tipicas en una orbita, y la potencia de cada uno. Se calcula la energia
por orbita.

\textbf{Ejemplo simplificado para tu payload}:
\[
E_{\text{orbita}} = \sum_i P_i \cdot t_i
\]

Tu reporte estima:
\begin{itemize}[leftmargin=2em]
  \item Idle: $3.0$ W $\times$ $80$ min $= 14400$ J
  \item 4 scans modo single: $23$ J $\times$ 4 $= 92$ J
  \item Comms: $4$ W $\times$ $10$ min $= 2400$ J
  \item \textbf{Total}: $\sim 17$ kJ por orbita $= 4.7$ Wh
\end{itemize}

\chapter{EPS satelital}

Un EPS (Electrical Power System) tiene 4 bloques principales:

\begin{itemize}[leftmargin=2em]
  \item \textbf{Paneles solares}: generan $5$--$30$ W segun tamaño y
        eficiencia.
  \item \textbf{MPPT} (Maximum Power Point Tracker): optimiza la extraccion.
  \item \textbf{Bateria}: $10$--$80$ Wh.
  \item \textbf{Distribucion}: rails de salida ($3.3$ V, $5$ V, $12$ V) con
        proteccion contra cortos.
\end{itemize}

% =====================================================================
\part{Aplicacion al CubeSat}

\chapter{El stack del payload INTISAT}

Resumen de todo lo que vimos aplicado a tu sistema:

\begin{itemize}[leftmargin=2em]
  \item \textbf{Energia}: $5$V\_PAYLOAD desde EPS $\to$ LDO $\to$ Pi 5 $\to$ PMIC interna $\to$ todos los rails internos.
  \item \textbf{Comunicacion}: I2C\_2 entre OBC y Pi via TXS0108E.
  \item \textbf{Sensor}: OV5647 por CSI-2 directo a la Pi.
  \item \textbf{Actuador}: OLED SSD1351 por SPI desde la Pi.
  \item \textbf{Monitor de potencia}: INA219 por I2C secundario (futuro).
\end{itemize}

\chapter{Diseño del carrier board (CM5)}

Cuando migren a CM5, el carrier board tiene que tener:

\begin{enumerate}[leftmargin=2em]
  \item Conector mezzanine de $200$ pines para el CM5 (Hirose DF40).
  \item LDO o switching regulator que tome $5$V\_PAYLOAD y lo regule a
        las tensiones que pide el CM5.
  \item Capacitores de decoupling de $100$ nF cerca de cada chip.
  \item Diodo Schottky de proteccion de polaridad inversa.
  \item Fusible PTC de $5$ A para corto-circuito.
  \item TXS0108E para I2C externo.
  \item INA219 para monitor de potencia.
  \item Conector PC-104 para integracion al INTISAT.
  \item Conector CSI-2 (15 o 22 pines) para OV5647.
  \item Conector SPI + GPIO para OLED.
  \item LEDs de status de boot.
  \item Pads de prueba para todos los buses.
\end{enumerate}

\chapter{Tests pre-vuelo}

Resumen de tests electronicos que tu equipo va a hacer:

\begin{itemize}[leftmargin=2em]
  \item \textbf{Continuidad y cortos}: con multimetro antes de encender.
  \item \textbf{Polaridad de fuente}: verificar antes de cada encendido.
  \item \textbf{Inrush current}: medir corriente de encendido inicial con osciloscopio.
  \item \textbf{Power-on sequence}: verificar que cada rail enciende en el orden correcto.
  \item \textbf{Burn-in}: 8 horas de operacion continua sin reboot.
  \item \textbf{TVAC}: ciclos termicos en vacio (-40 a +85 C).
  \item \textbf{Vibracion}: 3 g RMS, 5--2000 Hz, segun MIL-STD-810.
  \item \textbf{EMC}: emisiones radiadas y conducidas, susceptibilidad.
  \item \textbf{Funcional end-to-end}: 100 scans seguidos sin fallo.
\end{itemize}

% =====================================================================
\appendix
\chapter{Apendice --- Codigo de colores de resistores}

\begin{center}
\begin{tabular}{lcc}
\toprule
\textbf{Color} & \textbf{Digito / Multiplicador} & \textbf{Tolerancia} \\
\midrule
Negro    & 0 / $\times 1$        & --- \\
Marron   & 1 / $\times 10$       & $\pm 1\%$ \\
Rojo     & 2 / $\times 100$      & $\pm 2\%$ \\
Naranja  & 3 / $\times 1k$       & --- \\
Amarillo & 4 / $\times 10k$      & --- \\
Verde    & 5 / $\times 100k$     & $\pm 0.5\%$ \\
Azul     & 6 / $\times 1M$       & $\pm 0.25\%$ \\
Violeta  & 7 / $\times 10M$      & $\pm 0.1\%$ \\
Gris     & 8                      & --- \\
Blanco   & 9                      & --- \\
Dorado   & --- / $\times 0.1$    & $\pm 5\%$ \\
Plata    & --- / $\times 0.01$   & $\pm 10\%$ \\
\bottomrule
\end{tabular}
\end{center}

\chapter{Apendice --- Conversiones rapidas}

\begin{center}
\begin{tabular}{lll}
\toprule
\textbf{Concepto} & \textbf{Formula} & \textbf{Ejemplo} \\
\midrule
Corriente (A) a mA   & $\times 1000$    & $1.5$ A $= 1500$ mA \\
mAh a J              & $\times V \times 3.6$ & $1000$ mAh @ $5$ V $= 18\,000$ J \\
J a Wh               & $\div 3600$       & $7200$ J $= 2$ Wh \\
W a J/s              & $1:1$             & 5 W $= 5$ J/s \\
\bottomrule
\end{tabular}
\end{center}

\chapter{Apendice --- Hoja de ruta de aprendizaje}

\section{Plan de 3 meses para profundizar}

\textbf{Mes 1}: Capitulos 1--7. Construir un protoboard con resistencias,
capacitores, LEDs. Implementar circuitos basicos de Ohm y Kirchhoff.

\textbf{Mes 2}: Capitulos 8--13. Probar reguladores LDO y switching reales,
ver como cambia la temperatura. Programar un microcontrolador (Arduino o
ESP32).

\textbf{Mes 3}: Capitulos 14--20. Implementar I2C, SPI, UART en un proyecto
real. Diseñar y ordenar tu primer PCB en JLCPCB.

\section{Recursos para profundizar}

\begin{itemize}[leftmargin=2em]
  \item \textbf{Libro}: \emph{The Art of Electronics} (Horowitz \& Hill).
        El "Gonzalez \& Woods" de electronica.
  \item \textbf{Curso}: MIT 6.002 Circuits and Electronics (gratis online).
  \item \textbf{Practica}: TinkerCAD Circuits para simulacion online sin
        comprar nada.
  \item \textbf{Tutorial PCB}: KiCad oficial (gratis, profesional).
  \item \textbf{Comunidad}: \url{https://www.eevblog.com/forum/} y r/electronics.
\end{itemize}

\backmatter
\chapter*{Bibliografia}
\addcontentsline{toc}{chapter}{Bibliografia}

\begin{enumerate}[leftmargin=2em]
\item Horowitz P, Hill W. \emph{The Art of Electronics}. 3ed. Cambridge
  University Press, 2015.
\item Sedra A, Smith K. \emph{Microelectronic Circuits}. 8ed. Oxford, 2020.
\item Razavi B. \emph{Fundamentals of Microelectronics}. 2ed. Wiley, 2014.
\item Williams J. \emph{Analog Circuit Design}. Newnes, 1991.
\item Texas Instruments. \emph{INA219 Datasheet}. \url{https://www.ti.com/product/INA219}.
\item Raspberry Pi Foundation. \emph{BCM2712 Peripherals}. \url{https://datasheets.raspberrypi.com/bcm2712/}.
\item OmniVision. \emph{OV5647 Datasheet}, 2009.
\item Solomon Systech. \emph{SSD1351 Datasheet}, 2014.
\end{enumerate}

\end{document}
